% Image credits for the photographs and AI illustrations of Book 2 (High School), French edition.
% Machine-readable license records live in images/book2/CREDITS.md; keep this
% file in sync with frontmatter/image-credits-book2.tex (one \item per item there).
% Every credit, author, dataset and licence is kept exactly as in the English file.
\chapter*{Crédits des images}

Les dessins, cartes et graphiques de ce livre sont originaux~; les cartes
et graphiques tracés à partir de données publiées en citent la source dans
la légende ou dans cette liste. Les photographies sont utilisées sous leurs
licences libres respectives, avec nos remerciements à leurs auteurs~; les
liens complets vers les sources et les adresses des licences de chaque
photographie figurent dans \texttt{images/book2/CREDITS.md}, dans le dépôt
du livre.

\bigskip
À côté des photographies et des dessins originaux, certaines figures sont
des illustrations générées par IA, produites avec le générateur d'images
d'OpenAI à partir de consignes rédigées par les éditeurs du livre, et
vérifiées pour leur exactitude scientifique avant d'être retenues. Elles
ne montrent que des scènes génériques, jamais un lieu réel, une personne
réelle ou un spécimen réel. La consigne complète de chaque image générée
est consignée dans \texttt{images/book2/ai/PROMPTS.md}, dans le dépôt.

\bigskip
\noindent Photographies et données, par chapitre~:
\begin{itemize}\small
  \item Chap.~1 (mesurer la Terre)~: la Terre vue d'Apollo~17~; photographie NASA, domaine public (Wikimedia Commons).
  \item Chap.~1~: un cercle répétiteur, planche de J.-B. Delambre, \emph{Base du système métrique décimal}, vol.~2 (1807)~; domaine public (Wikimedia Commons, d'après BnF Gallica).
  \item Chap.~1~: longueurs du degré calculées sur l'ellipsoïde WGS~84 (jeu de données EPSG).
  \item Chap.~2 (une planète habitable)~: Vénus vue par Mariner~10 (PIA23791), NASA/JPL-Caltech, domaine public~; le mont Sharp vu par le rover Curiosity (PIA16769), NASA/JPL-Caltech/MSSS, domaine public (Wikimedia Commons).
  \item Chap.~2~: données planétaires des fiches de la NASA (NASA fact sheets)~; limites des phases de l'eau d'après les publications de l'IAPWS SR1-86 et R14-08~; limites de la zone habitable d'après Kopparapu et ses collègues (2013)~; Proxima~b d'après la NASA Exoplanet Archive.
  \item Chap.~3 (le bilan énergétique de la Terre)~: mares de fonte sur la banquise arctique (mission ICESCAPE, 2011)~; photographie du NASA Goddard Space Flight Center, domaine public (Wikimedia Commons).
  \item Chap.~3~: valeurs du bilan énergétique d'après le sixième rapport d'évaluation du GIEC (IPCC Sixth Assessment Report), groupe de travail~I, figure~7.2~; albédos d'après l'University of Colorado Boulder (MOSAiC) et le National Snow and Ice Data Center~; transport de chaleur d'après Trenberth et Caron (2001)~; production primaire nette d'après Field et ses collègues (1998)~; données thermochimiques du NIST Chemistry WebBook.
  \item Chap.~4 (pourquoi les nuages se forment)~: cumulus de beau temps, photographie de Ralph F. Kresge, NOAA~; l'enclume d'un cumulonimbus, NOAA National Severe Storms Laboratory~; tous deux domaine public (Wikimedia Commons).
  \item Chap.~4~: saturation calculée d'après les publications de l'IAPWS~; tailles des gouttes d'après le NASA Jet Propulsion Laboratory.
  \item Chap.~5 (orages et cyclones tropicaux)~: l'ouragan Isabel vu depuis la Station spatiale internationale, photographie de Mike Trenchard, NASA Johnson Space Center~; une tornade au Kansas, photographie de Sean Waugh, NOAA National Severe Storms Laboratory~; des éclairs, photographie de C.~Clark, NOAA National Severe Storms Laboratory~; tous domaine public (Wikimedia Commons).
  \item Chap.~5~: trajectoires des cyclones d'après IBTrACS version~4r01 (Gahtan et ses collègues, 2024~; Knapp et ses collègues, 2010~; NOAA National Centers for Environmental Information)~; catégories d'ouragans d'après le NOAA National Hurricane Center~; énergie et conditions de formation des ouragans d'après le NOAA Atlantic Oceanographic and Meteorological Laboratory~; fréquence des éclairs d'après Christian et ses collègues (2003)~; données sur les tornades d'après le NOAA Storm Prediction Center.
  \item Chap.~6 (un océan en couches)~: une rosette de prélèvement, photographie de Pablo Clemente-Colon, NOAA Photo Library~; un flotteur profileur mis à l'eau depuis un navire de recherche, NOAA Global Ocean Monitoring and Observing Program~; la chlorophylle océanique 2002--2010, Jesse Allen et Robert Simmon, NASA Earth Observatory~; tous domaine public (Wikimedia Commons).
  \item Chap.~6~: profils océaniques d'après le World Ocean Atlas 2023 (Reagan et ses collègues, 2024~; NOAA National Centers for Environmental Information)~; masse volumique de l'eau de mer d'après l'équation d'état internationale (Fofonoff et Millard, 1983)~; volume de l'océan d'après Charette et Smith (2010)~; données sur les flotteurs d'après l'Argo Program Office~; remontées d'eau d'après le NOAA National Ocean Service et Kessler (2006).
  \item Chap.~7 (le sol, une ressource lente à se former)~: un profil de latérite sur gneiss au Kerala, photographie de Werner Schellmann, CC~BY-SA~2.5~; une saprolite du centre de Madagascar, photographie de Diorit, CC~BY-SA~3.0~; un profil de sol, USDA Natural Resources Conservation Service, domaine public (tous Wikimedia Commons).
  \item Chap.~7~: vitesses de production et d'érosion des sols d'après Montgomery (2007)~; formules des minéraux d'après la Mineralogy Database (webmineral.com)~; épaisseur des saprolites d'après le USGS Ground Water Atlas of the United States.
  \item Chap.~8 (bassins versants, rivières et crues)~: un hydrologue mesurant un débit et une station de jaugeage, U.S. Geological Survey, domaine public (Wikimedia Commons).
  \item Chap.~8~: pointes annuelles et débits journaliers du Potomac près de Washington, U.S. Geological Survey National Water Information System, station 01646500~; débit et bassin du Mississippi d'après Kammerer (1990, USGS)~; débit de l'Amazone d'après Callède et ses collègues (2010)~; sécurité face aux eaux de crue d'après le NOAA National Weather Service.
  \item Chap.~9 (les ressources en eau douce)~: la mer d'Aral en 2000 et en 2018, NASA Earth Observatory (MODIS/Terra)~; champs irrigués près de Sublette, au Kansas, NASA (ASTER/Terra), domaine public (Wikimedia Commons).
  \item Chap.~9~: stocks d'eau mondiaux d'après Shiklomanov (1993), U.S. Geological Survey Water Science School~; cycle de l'eau mondial d'après Trenberth et ses collègues (2007), tel qu'il est publié dans Trenberth et Asrar (2014)~; âges des eaux souterraines nubiennes d'après Sturchio et ses collègues (2004)~; variations de l'aquifère des High Plains d'après McGuire (2017, USGS)~; bilan de la mer d'Aral d'après Gaybullaev et ses collègues (2012) et Micklin (2007).
  \item Chap.~10 (de la dérive des continents à l'expansion des fonds océaniques)~: le navire de forage \emph{JOIDES Resolution}, Ocean Drilling Program, Texas A\&M University, domaine public (Wikimedia Commons).
  \item Chap.~10~: l'échelle des inversions de polarité de Cande et Kent (1995)~; les résultats de forage du Leg~3 du Deep Sea Drilling Project (1970)~; l'âge des fonds océaniques d'après Seton et ses collègues (2020), EarthByte, licence CC~BY~4.0.
  \item Chap.~11 (les ondes sismiques révèlent l'intérieur)~: installation d'un sismomètre temporaire à Porto Rico~; photographie de Thomas Pratt, U.S. Geological Survey, domaine public (Wikimedia Commons).
  \item Chap.~11~: le Preliminary Reference Earth Model de Dziewonski et Anderson (1981)~; temps de trajet empiriques des ondes P et S d'après l'Australian Passive Seismic Server (AuSPASS).
  \item Chap.~12 (lithosphère et asthénosphère)~: un iceberg, NOAA, domaine public (USGS Water Science School).
  \item Chap.~12~: flux géothermique d'après Davies et Davies (2010), licence CC~BY~3.0~; provinces de flux géothermique d'après Roy, Blackwell et Decker (1972)~; le solidus du manteau de Hirschmann (2000)~; épaisseur et masse volumique de la croûte d'après Christensen et Mooney (1995)~; altitudes d'après ETOPO1 (NOAA).
  \item Chap.~13 (les plaques et leurs mouvements)~: une station de positionnement par satellite dans la chaîne des Cascades~; photographie de Dave Sherrod, U.S. Geological Survey, domaine public (Wikimedia Commons). Relief ombré du fond du Pacifique Nord, National Geophysical Data Center, NOAA, domaine public (Wikimedia Commons).
  \item Chap.~13~: pôles de rotation des plaques NNR-MORVEL56 d'Argus, Gordon et DeMets (2011)~; positions des stations d'après le Nevada Geodetic Laboratory (Blewitt, Hammond et Kreemer, 2018)~; âges des volcans de Clague et Dalrymple (1987), U.S. Geological Survey~; limites de plaques de Bird (2003)~; terres d'après Natural Earth.
  \item Chap.~14 (les dorsales océaniques et la lithosphère océanique)~: laves en coussins sur le rift des Galápagos, NOAA Okeanos Explorer Program, domaine public~; un fumeur noir sur la dorsale Juan de Fuca, University of Washington et NOAA, domaine public~; gabbro, photographie de Darla Sondrol (GeoDIL), CC0~; péridotite dans un basalte, photographie de James St.~John, CC~BY~2.0 (tous Wikimedia Commons).
  \item Chap.~14~: profondeurs des fonds océaniques d'après ETOPO1 (NOAA) et âges d'après Seton et ses collègues (2020), EarthByte, licence CC~BY~4.0~; lois de profondeur de Parsons et Sclater (1977)~; le solidus du manteau de Hirschmann (2000)~; pôles de rotation NNR-MORVEL56 d'Argus, Gordon et DeMets (2011).
  \item Chap.~15 (les zones de subduction)~: schiste bleu et éclogite, photographies de James St.~John, CC~BY~2.0~; dégâts du tsunami à Kesennuma, au Japon, en 2011, photographie de Dylan McCord, U.S. Navy, domaine public (tous Wikimedia Commons).
  \item Chap.~15~: séismes d'après le catalogue de séismes de l'USGS (ComCat)~; limites de plaques et vitesses de convergence de Bird (2003)~; âges des fonds océaniques d'après Seton et ses collègues (2020), EarthByte, licence CC~BY~4.0~; vitesses des plaques NNR-MORVEL56 d'Argus, Gordon et DeMets (2011)~; terres d'après Natural Earth.
  \item Chap.~16 (les minéraux)~: quartz, orthose, biotite et olivine, photographies de Shannon Heinle, Darla Sondrol, Dexter Perkins et Nessa Eull (Geology Digital Image Library), CC0~; basalte en lame mince, photographies de Chmee2, CC~BY-SA~3.0~; quartz d'un granite en lame mince, photographie de Michael C.~Rygel, CC~BY-SA~3.0 (tous Wikimedia Commons).
  \item Chap.~16~: la composition de la croûte continentale de Rudnick et Gao (2014)~; masses atomiques standard de l'IUPAC/CIAAW~; formules et masses volumiques des minéraux d'après webmineral.com~; masses volumiques des roches granitiques de la Sierra Nevada d'après le U.S. Geological Survey (Open-File Report 91-570-A).
  \item Chap.~17 (les roches magmatiques)~: colonnes de basalte de la Chaussée des Géants, photographie d'Adam Bishop, CC~BY-SA~3.0~; El Capitan, Yosemite, U.S. Geological Survey, CC0~; basalte, andésite, rhyolite, diorite, granite, obsidienne, une roche porphyrique et deux lames minces, photographies de Nessa Eull, Twyla Baker, Joseph H.~Hartman, Mike Davis, Dexter Perkins et Darla Sondrol (Geology Digital Image Library), CC0 (tous Wikimedia Commons).
  \item Chap.~17~: refroidissement des dykes et des sills d'après Delaney (1986), U.S. Geological Survey Open-File Report 86-444~; compositions des magmas et ordre de cristallisation d'après les notes de géologie physique de S.~A.~Nelson et S.~Earle, \emph{Physical Geology} (CC~BY~4.0)~; biographie de Bowen d'après l'histoire du Geophysical Laboratory publiée par Carnegie Science.
  \item Chap.~18 (la géologie du charbon, du pétrole et du gaz)~: plantes fossiles dans une argilite riche en carbone, lignite, charbon bitumineux et anthracite, photographies de James St.~John, CC~BY~2.0~; une couche de lignite en Roumanie, U.S. Geological Survey, domaine public (tous Wikimedia Commons).
  \item Chap.~18~: fenêtre à huile, formation du charbon et pièges d'après S.~Earle, \emph{Physical Geology} (CC~BY~4.0)~; teneur en carbone du charbon d'après la U.S. Energy Information Administration~; le système pétrolier d'après Magoon et Schmoker (U.S. Geological Survey, 2000)~; les bassins de rift et les bassins salifères de l'ouest de l'Afrique centrale d'après la U.S. Geological Survey Fact Sheet 2023-3028.
  \item Chap.~19 (la croûte continentale)~: le gneiss d'Acasta, photographie de Pedroalexandrade, CC~BY-SA~3.0~; zircons des Jack Hills, images d'A.~Cavosie et J.~Valley, domaine public~; gneiss et granites du Bouclier canadien, photographie de Chris Light, CC~BY-SA~4.0 (tous Wikimedia Commons).
  \item Chap.~19~: composition de la croûte de Rudnick et Gao (2014)~; épaisseur de la croûte d'après Christensen et Mooney (1995), White, McKenzie et O'Nions (1992) et Karplus et ses collègues (2019)~; étendue des continents d'après Cogley (1984)~; âges des fonds océaniques d'après Seton et ses collègues (2020), EarthByte, licence CC~BY~4.0.
  \item Chap.~20 (la formation des montagnes)~: l'Everest vu de la navette spatiale Columbia, NASA, domaine public~; couches de roches plissées dans les gorges de la Kali Gandaki, photographie de Jean-Marie Hullot, CC~BY~2.0~; le chevauchement de Glaris, photographie de Christian Heine, CC~BY-SA~3.0~; laves en coussins de l'ophiolite de Semail, photographie de Lionel Duchoiselle, CC~BY-SA~3.0 (tous Wikimedia Commons).
  \item Chap.~20~: positions passées de l'Inde et de Lhassa d'après les modèles de plaques de Müller et ses collègues (2019, 2022), calculées avec le GPlates Web Service (EarthByte)~; reliefs élevés d'après ETOPO1~; séismes d'après le catalogue de l'USGS~; limites de plaques de Bird (2003)~; terres d'après Natural Earth.
  \item Chap.~21 (les roches métamorphiques)~: schistosité ardoisière au microscope, photographie de Geoloshe, CC0~; micaschiste à grenat de Syros, photographie de Graeme Churchard, CC~BY~2.0~; schiste à disthène et gneiss à quartz et biotite, photographies de James St.~John, CC~BY~2.0 (tous Wikimedia Commons).
  \item Chap.~21~: les domaines de stabilité du disthène, de l'andalousite et de la sillimanite d'après Holdaway et Mukhopadhyay (1993).
  \item Chap.~22 (le magmatisme des zones de subduction)~: le panache de l'éruption du mont Saint Helens, 18~mai 1980, photographie de Donald A. Swanson, USGS, domaine public~; granodiorite, andésite et Sierra Nevada, photographies de James St.~John, CC~BY~2.0 (tous Wikimedia Commons).
  \item Chap.~22~: séismes d'après le catalogue de l'USGS~; volcans d'après le Global Volcanism Program de la Smithsonian~; limites de plaques de Bird (2003)~; terres d'après Natural Earth~; le solidus de la péridotite sèche de Hirschmann (2000), avec les masses volumiques de PREM.
  \item Chap.~23 (la destruction des montagnes)~: les Appalaches vues de la Station spatiale internationale, photographie d'astronaute ISS033-E-22378, NASA, domaine public~; conglomérat à galets de quartz, photographie de James St.~John, CC~BY~2.0 (tous deux Wikimedia Commons).
  \item Chap.~23~: taux de dénudation de Portenga et Bierman (2011)~; charges sédimentaires du Mississippi de Meade et Moody (2010)~; constante de décroissance d'Ahnert (1970)~; les collines de Varden, la région imaginaire de la série.
  \item Chap.~24 (la chaleur interne de la Terre)~: la centrale géothermique de Nesjavellir, en Islande, photographie de Gretar Ívarsson, retouchée par Fir0002, domaine public (Wikimedia Commons).
  \item Chap.~24~: flux géothermique de Davies et Davies (2010)~; chaleur radiogénique d'après la mesure des géoneutrinos de KamLAND (2011) et la chaleur libérée aujourd'hui par les quatre isotopes~; production de chaleur des roches d'après Vilà et Fernández (2010).
  \item Chap.~25 (la datation radiométrique et l'âge de la Terre)~: William Thomson, lord Kelvin, en 1906, photographie de Russell and Sons, Londres, Smithsonian Institution Libraries, sans restriction de droits connue~; la météorite de fer de Canyon Diablo, photographie de Geoffrey Notkin, CC~BY-SA~2.5 (tous deux Wikimedia Commons).
  \item Chap.~25~: isotopes du plomb des météorites de Patterson (1956)~; la constante de désintégration du rubidium~87 recommandée par l'IUGS et l'IUPAC (2022)~; l'histoire des estimations, d'après Holmes (1913, 1927), Patterson (1956) et l'U.S. Geological Survey~; les collines de Varden, la région imaginaire de la série.
  \item Chap.~26 (la Terre primitive et son atmosphère)~: formation de fer rubané, Dales Gorge, photographie de Graeme Churchard, CC~BY~2.0~; stromatolithes à Shark Bay, photographie de Paul Harrison, CC~BY-SA~3.0~; couches rouges près de Sedona, photographie de Wendy from Pennsylvania, CC~BY-SA~2.0 (tous Wikimedia Commons).
  \item Chap.~26~: la luminosité du jeune Soleil d'après Gough (1981)~; limites et couleurs des éons d'après la Charte chronostratigraphique internationale.
  \item Chap.~27 (lire les climats du passé)~: bulles d'air dans la glace antarctique, photographie CSIRO Atmospheric Research, CC~BY~3.0~; foraminifères planctoniques, photographie de Marci Robinson, USGS, domaine public~; grains de pollen, Dartmouth Electron Microscope Facility, domaine public (tous Wikimedia Commons).
  \item Chap.~27~: température EPICA Dôme~C (Jouzel et ses collègues, 2007) et dioxyde de carbone des carottes de glace antarctiques (Bereiter et ses collègues, 2015) via NOAA Paleoclimatology~; l'empilement LR04 (Lisiecki et Raymo, 2005)~; relief ETOPO1 (NOAA)~; NASA GISS (GISTEMP v4)~; périodes de référence d'après le sixième rapport d'évaluation du GIEC (IPCC Sixth Assessment Report), groupe de travail~I, chapitre~2.
  \item Chap.~28 (le changement climatique)~: l'observatoire du Mauna Loa~; photographie de Brian Vasel, NOAA, domaine public (Wikimedia Commons).
  \item Chap.~28~: Charles David Keeling et John Chin au Mauna Loa, 1997~; photographie de John Miller, NOAA, domaine public (Wikimedia Commons).
  \item Chap.~28~: données du NOAA Global Monitoring Laboratory et de la Scripps Institution of Oceanography (dioxyde de carbone), de la NASA GISS (GISTEMP v4), du NOAA Laboratory for Satellite Altimetry (niveau de la mer), de Bereiter et ses collègues (2015) via NOAA Paleoclimatology (carottes de glace), et du sixième rapport d'évaluation du GIEC (IPCC Sixth Assessment Report), groupe de travail~I (forçages, attribution, projections).
\end{itemize}
